\section{UV Divergences at One Loop}
\label{sec:UVOL}
We present the relevant one-loop Feynman diagrams for quasi-gluon operators of an asymptotic gluon of momentum $p$ in Fig.~\ref{fig:oneloop-g}, where the bubble in the diagram (e) includes all one-loop self-energy diagrams of the active gluon.  For the complete one-loop contribution, additional Feynman diagrams are needed.  Some of them are mirror of diagrams (b), (c), (d), (e) and (g), while the rest can be obtained by replacing external momentum $p$ to $- p$ in all these Feynman amplitudes.  For the following one-loop calculation, we take the linearly combined quasi-gluon operator in Ref.~\cite{Ji:2013dva} as an example, but our conclusion is true for any of the 36 independent operators.

%%%%%%%%%%%%%%%%%%%%%%%%%%%%%%%
\begin{figure}[htb]
\begin{center}
\includegraphics[width=0.75\textwidth]{images/oneloop-g.pdf}
\caption{\label{fig:oneloop-g}
Some typical Feynman diagrams for quasi-gluon PDFs of an asymptotic gluon of momentum $p$ at one-loop order.}
\end{center}
\end{figure}
%%%%%%%%%%%%%%%%%%%%%%%%%%%%%%

We choose Feynman gauge, and assume $\xi_z$ to be positive for definiteness.  The diagram (a) in Fig.~\ref{fig:oneloop-g} gives
\begin{align}
M_{1{a}}&= \frac{g_s^2 \mu_r^{4-d} C_A}{i} \,e^{-i p_z \xi_z} \int_0^{\xi_z}dr_1\int_{r_1}^{\xi_z}dr_2 \frac{d^d l}{(2\pi)^d} \frac{e^{il_z(r_2-r_1)}}{l^2}
\nonumber\\
&\xlongequal{\text{UV}} \frac{\alpha_s C_A}{\pi} e^{-i p_z \xi_z} \left(-\frac{\pi \mu_r\xi_z}{3-d}+ \frac{1}{4-d} \right) \, ,
\label{eq:1a0}
\end{align}
where $\mu_r$ is a renormalization scale to compensate the mass dimension in DR. This diagram contributes to both linear and logarithmic UV divergences, as expected.

To understand where in this one-loop phase space the UV divergences in Eq.~(\ref{eq:1a0}) come from, it is instructive to distinguish $l_z$ - the $z$-component of the loop momentum $l$ from $\bar{l}_\mu$ - the other components of $l$, as $l_\mu = \bar{l}_\mu - l_z n_{\mu}$ with $l^2=\bar{l}^2-l_z^2$ \cite{Ishikawa:2017faj}.
If $\bar{l}^2$ is constrained in a finite region in Eq.~(\ref{eq:1a0}), integrating $l_z$, $r_1$ and $r_2$ cannot generate any UV divergence.  Furthermore, there is no UV divergence if we do not include the region where $|r_2-r_1|$ is very small, which can be demonstrated by introducing the following decomposition,
\begin{align}\label{eq:dec}
\frac{1}{(l+q)^2}
= \frac{1}{(\bar{l}+\bar{q})^2}+ \frac{(l_z+q_z)^2}{(l+q)^2(\bar{l}+\bar{q})^2}\, ,
\end{align}
where $q$ can be any unintegrated momentum. By applying this decomposition to $\frac{1}{l^2} $ in Eq.~\eqref{eq:1a0}, the first term is free of $l_z$, and thus the integration of $l_z$ gives $\delta(r_2-r_1)$, while the second term is UV finite under the integration of $\bar{l}$.  That is, the UV divergence in Eq.~\eqref{eq:1a0} can only come from the region of phase space where $\bar{l}$ are in UV region while $|r_2-r_1|$ is very small. Therefore, we conclude that, with DR, all UV divergences of diagram (a) in Fig.~\ref{fig:oneloop-g} come from a region localized in spacetime.

By decomposing both $\frac{1}{l^2} $ and $\frac{1}{(p-l)^2}$ using Eq.~\eqref{eq:dec}, we obtain many terms for diagrams (b) and (c) and found that these terms are either free of $l_z$ in denominator, which result in $\delta(r)$ or its derivatives, or UV finite under the integration of $\bar{l}$. Thus the UV divergences of these two diagrams are also localized in spacetime,
\begin{align}
M_{1{b}}
&\xlongequal{\text{UV}} \frac{\alpha_s C_A}{\pi} e^{-i p_z \xi_z} \left(\frac{-i}{p_z \xi_z} \frac{\pi \mu_r \xi_z}{3-d} \right) \, , \label{eq:1b0}\\
M_{1c}
&\xlongequal{\text{UV}} \frac{\alpha_s C_A}{\pi} e^{-i p_z \xi_z} \left(\frac{i}{p_z \xi_z} \frac{\pi \mu_r \xi_z}{3-d}+\frac{3}{4} \frac{1}{4-d} \right) \, ,
\label{eq:1c0}
\end{align}
where diagram (b) has only linear UV divergence, while (c) has both linear and logarithmic UV divergence.

The diagrams (d) and (e) in Fig.~\ref{fig:oneloop-g} have only logarithmic UV divergences, which come from the region where all components of $l^\mu$ go to infinity, and thus, is localized in spacetime.
All other diagrams in Fig.~\ref{fig:oneloop-g} are free of UV divergence, simply because the loop cannot be localized in spacetime due to finite $\xi_z$, same as the argument for quasi-quark operators \cite{Ishikawa:2017faj}.
In summary, we conclude that the UV divergences of quasi-gluon operators at one-loop can only be emerged from a region {\it localized}\ in coordinate spacetime. For a comparison, UV divergences of operators defining PDFs come from a region   {\it nonlocal}\ along the ``$-$" light-cone direction \cite{Ishikawa:2017faj}.

The linear UV divergence in Eq.~(\ref{eq:1a0}) from the diagram (a) is harmless because it can be easily exponentiated to all-order as an overall phase factor and then multiplicatively renormalized, just like the case of quasi-quark operators \cite{Ishikawa:2017faj,Polyakov:1980ca,Dotsenko:1979wb}. However, the presence of linear UV divergence in Eqs.~(\ref{eq:1b0}) and (\ref{eq:1c0}) from the diagrams (b) and (c), respectively, could challenge the multiplicative renormalizability. Fortunately, we find that with DR, the linear UV divergences from these two diagrams are canceled.  On the contrary, the linear divergences from the diagrams (b) and (c) do not cancel if one uses a cutoff regularization that breaks the gauge symmetry \cite{Wang:2017qyg,Wang:2017eel}. This implies that gauge invariance plays an important role to remove the linear UV divergences that may challenge the multiplicative renormalizability.  In the following, we will use gauge invariance to show that all linear divergences, except that from self-energy of gauge links, are canceled by summing over all contributions, and quasi-gluon operators could be multiplicatively renormalized.
